% Compile with: pdflatex RIS_MAIN_BEAM_AND_SIDELOBE_DITHERING.tex
\documentclass[11pt,a4paper]{article}
\usepackage{amsmath,amssymb,bm,geometry}
\usepackage[hidelinks]{hyperref}
\geometry{margin=2.2cm}
\newcommand{\C}{\mathbb{C}}
\newcommand{\rx}{\mathrm{rx}}
\title{RIS: Main beam toward RX + sidelobe-only environment scan}
\author{}
\date{}
\begin{document}
\maketitle

This note reformulates the idea of \textbf{keeping the primary (focal) response aligned to the receiver} while using \textbf{dithering} to probe the environment---conceptually scanning \textbf{sidelobes} or ambient radiation---without destroying the main link.

\section{What you are trying to do}

You want a \textbf{two-part} surface modulation on the RIS:

\begin{center}
\begin{tabular}{|l|p{11cm}|}
\hline
\textbf{Main term} & Keeps the RIS-assisted contribution \textbf{strong and stable} at the RX (the ``beam'' or \textbf{focus} pointed at the RX). \\
\hline
\textbf{Dither / scan term} & Changes how energy is distributed \textbf{elsewhere} (sidelobes / off-focus regions) to \textbf{sense} scatterers, blockers, or changes---ideally with \textbf{minimal} impact on the RX. \\
\hline
\end{tabular}
\end{center}

Important: on a finite aperture you cannot literally move ``only'' sidelobes with \textbf{zero} effect on the main coupling unless you \textbf{constrain} how you dither. The rigorous formulation is: \textbf{keep the perturbation small and/or orthogonal to the RX coupling}.

\section{Discrete RIS model (schematic)}

Model the RIS by $N$ unit cells. Let the \textbf{complex spatial modulation} (per cell) be stacked in a vector:
\begin{equation}
\bm{\Gamma} = [\Gamma_1,\, \Gamma_2,\, \ldots,\, \Gamma_N]^{\mathsf{T}} \in \C^N.
\end{equation}

The \textbf{scalar} (or dominant) complex gain seen at the RX through the RIS path can be written schematically as an \textbf{inner product} with a \textbf{coupling vector} from the elements to the RX:
\begin{equation}
h_{\mathrm{RIS}} \propto \mathbf{g}_{\rx}^{\mathsf{H}}\, \bm{\Gamma}.
\end{equation}

\begin{itemize}
\item $\mathbf{g}_{\rx} \in \C^N$ encodes how each element couples to the RX (phases, amplitudes, path aggregation).
\item Your \textbf{baseline} ``point the main toward RX'' pattern is some $\bm{\Gamma}_0$ chosen so that $\mathbf{g}_{\rx}^{\mathsf{H}} \bm{\Gamma}_0$ is large (after any normalization you use).
\end{itemize}

This is the right abstraction for both \textbf{far-field steering} (where $\mathbf{g}_{\rx}$ looks like a plane-wave array response) and \textbf{near-field focusing} (where $\mathbf{g}_{\rx}$ includes \textbf{range and curvature}, not just angle).

\section{Splitting main and dither}

\subsection{Additive complex perturbation}
\begin{equation}
\bm{\Gamma}_k = \bm{\Gamma}_0 + \Delta\bm{\Gamma}_k.
\end{equation}
\begin{itemize}
\item $\bm{\Gamma}_0$: fixed (or slowly tracked) \textbf{main} pattern toward the RX.
\item $\Delta\bm{\Gamma}_k$: $k$-th \textbf{dither} pattern for environment scanning.
\end{itemize}

You usually \textbf{renormalize} $\bm{\Gamma}_k$ to satisfy hardware constraints (e.g.\ per-element magnitude limits).

\subsection{Phase-only dither (common for passive RIS)}
\begin{equation}
\Gamma_{k,n} = \Gamma_{0,n}\, e^{j \delta_{k,n}}, \qquad n = 1,\ldots,N.
\end{equation}
\begin{itemize}
\item $\bm{\delta}_k = [\delta_{k,1},\ldots,\delta_{k,N}]^{\mathsf{T}}$: small or structured \textbf{phase ripple} on top of the main profile.
\end{itemize}

\section{Two ways to get ``main stays, environment is probed''}

\subsection{Null-space / orthogonal dither (strongest interpretation)}

You want $\Delta\bm{\Gamma}_k$ to \textbf{not} move the RX coupling, in first order:
\begin{equation}
\mathbf{g}_{\rx}^{\mathsf{H}}\, \Delta\bm{\Gamma}_k \approx 0.
\end{equation}

So $\Delta\bm{\Gamma}_k$ lies approximately in the \textbf{orthogonal complement} of $\mathbf{g}_{\rx}$.

\textbf{Construction (conceptual):}
\begin{enumerate}
\item Form $\mathbf{g}_{\rx}$ (from geometry + your near-field or far-field model, or from a simulator).
\item Build an orthonormal basis $\{\mathbf{u}_1, \ldots, \mathbf{u}_M\}$ orthogonal to $\mathbf{g}_{\rx}$:
\begin{equation}
\mathbf{g}_{\rx}^{\mathsf{H}} \mathbf{u}_m = 0, \quad m = 1,\ldots,M.
\end{equation}
\item Define dither codewords as
\begin{equation}
\Delta\bm{\Gamma}_k = \sum_{m=1}^{M} \alpha_{k,m}\, \mathbf{u}_m,
\end{equation}
with small coefficients $\alpha_{k,m}$.
\item Full pattern:
\begin{equation}
\bm{\Gamma}_k = \mathrm{Normalize}\bigl(\bm{\Gamma}_0 + \Delta\bm{\Gamma}_k\bigr).
\end{equation}
\end{enumerate}

\textbf{Intuition:} the RX inner product stays dominated by $\bm{\Gamma}_0$; the dither mostly changes \textbf{how energy is distributed away from} the RX coupling---what you loosely call \textbf{sidelobes} or non-focal radiation.

\subsection{Small high--spatial-frequency ripple (simpler, less controlled)}

Add a \textbf{weak}, \textbf{rapidly varying} phase $\bm{\delta}_k$ on top of $\bm{\Gamma}_0$. Often leaves the coherent sum toward RX mostly unchanged if ripple is small and balanced; still perturbs off-boresight structure. Easier to implement than null-space projection but typically causes \textbf{some} main-link variation.

\section{What the near-field case changes}

\begin{itemize}
\item \textbf{Far-field:} $\mathbf{g}_{\rx}$ is often well approximated by a \textbf{plane-wave} array steering vector.
\item \textbf{Near-field (Fresnel region):} $\mathbf{g}_{\rx}$ must reflect \textbf{distance and wavefront curvature} across elements---not only angle.
\end{itemize}

So \textbf{main toward RX} should be \textbf{near-field-aware} $\bm{\Gamma}_0$, and \textbf{null-space dither} should be orthogonal to this \textbf{near-field coupling vector}, not only to a far-field steering vector.

\section{Practical recipe}

\begin{enumerate}
\item Define $\mathbf{g}_{\rx}$ for your layout.
\item Set $\bm{\Gamma}_0$ to maximize $|\mathbf{g}_{\rx}^{\mathsf{H}} \bm{\Gamma}_0|$ (often $\bm{\Gamma}_0 \propto \mathbf{g}_{\rx}^{*}$ up to normalization).
\item Build basis $\{\mathbf{u}_m\}$ with $\mathbf{g}_{\rx}^{\mathsf{H}} \mathbf{u}_m = 0$.
\item Form $\bm{\Gamma}_k = \mathrm{Normalize}(\bm{\Gamma}_0 + \sum_m \alpha_{k,m} \mathbf{u}_m)$ with small $\alpha_{k,m}$.
\item Verify $|\mathbf{g}_{\rx}^{\mathsf{H}} \bm{\Gamma}_k|$ is stable across $k$.
\end{enumerate}

\section{Caveats}

\begin{enumerate}
\item \textbf{Exact} ``no change on main'' only holds in a \textbf{linear} sense; large phase-only dither still couples nonlinearly through $\exp(j\delta)$.
\item \textbf{Normalization} can couple dither and main when projecting to feasible $\bm{\Gamma}_k$.
\item In Sionna-style simulators you implement $\bm{\Gamma}_k$ as your own profiles on the RIS grid; this math is \textbf{your codebook design}.
\end{enumerate}

\section{One-line summary}

\textbf{Keep a near-field-correct main pattern $\bm{\Gamma}_0$ aligned to the RX; add dither $\Delta\bm{\Gamma}_k$ that is approximately orthogonal to $\mathbf{g}_{\rx}$ so the primary link is preserved while the rest of the field is modulated for environment sensing.}

\end{document}
